% Options for packages loaded elsewhere
\PassOptionsToPackage{unicode}{hyperref}
\PassOptionsToPackage{hyphens}{url}
\PassOptionsToPackage{dvipsnames,svgnames*,x11names*}{xcolor}
\PassOptionsToPackage{space}{xeCJK}
%
\documentclass[
  english,
  a4paper,
  DIV=12,
  footlines=2.1,
  usegeometry=true]{scrartcl}
\usepackage{amsmath,amssymb}
\usepackage{lmodern}
\usepackage{setspace}
\setstretch{1.1}
\usepackage{iftex}
\ifPDFTeX
  \usepackage[T1]{fontenc}
  \usepackage[utf8]{inputenc}
  \usepackage{textcomp} % provide euro and other symbols
\else % if luatex or xetex
  \usepackage{unicode-math}
  \defaultfontfeatures{Scale=MatchLowercase}
  \defaultfontfeatures[\rmfamily]{Ligatures=TeX,Scale=1}
  \setmainfont[]{Charis SIL}
  \setsansfont[]{Charis SIL Compact}
  \ifXeTeX
    \usepackage{xeCJK}
    \setCJKmainfont[]{Noto Serif CJK TC}
  \fi
  \ifLuaTeX
    \usepackage[]{luatexja-fontspec}
    \setmainjfont[]{Noto Serif CJK TC}
  \fi
\fi
% Use upquote if available, for straight quotes in verbatim environments
\IfFileExists{upquote.sty}{\usepackage{upquote}}{}
\IfFileExists{microtype.sty}{% use microtype if available
  \usepackage[]{microtype}
  \UseMicrotypeSet[protrusion]{basicmath} % disable protrusion for tt fonts
}{}
\usepackage{xcolor}
\IfFileExists{xurl.sty}{\usepackage{xurl}}{} % add URL line breaks if available
\IfFileExists{bookmark.sty}{\usepackage{bookmark}}{\usepackage{hyperref}}
\hypersetup{
  pdftitle={Horrordatei zum Testen vom neuen Textsatz},
  pdflang={en},
  pdfkeywords={Entangled Religions, religious contact, Markdown},
  colorlinks=true,
  linkcolor={NavyBlue},
  filecolor={Maroon},
  citecolor={Blue},
  urlcolor={Blue},
  pdfcreator={LaTeX via pandoc}}
\urlstyle{same} % disable monospaced font for URLs
\usepackage{longtable,booktabs,array}
\usepackage{calc} % for calculating minipage widths
% Correct order of tables after \paragraph or \subparagraph
\usepackage{etoolbox}
\makeatletter
\patchcmd\longtable{\par}{\if@noskipsec\mbox{}\fi\par}{}{}
\makeatother
% Allow footnotes in longtable head/foot
\IfFileExists{footnotehyper.sty}{\usepackage{footnotehyper}}{\usepackage{footnote}}
\makesavenoteenv{longtable}
\usepackage{graphicx}
\makeatletter
\def\maxwidth{\ifdim\Gin@nat@width>\linewidth\linewidth\else\Gin@nat@width\fi}
\def\maxheight{\ifdim\Gin@nat@height>\textheight\textheight\else\Gin@nat@height\fi}
\makeatother
% Scale images if necessary, so that they will not overflow the page
% margins by default, and it is still possible to overwrite the defaults
% using explicit options in \includegraphics[width, height, ...]{}
\setkeys{Gin}{width=\maxwidth,height=\maxheight,keepaspectratio}
\setlength{\emergencystretch}{3em} % prevent overfull lines
\providecommand{\tightlist}{%
  \setlength{\itemsep}{0pt}\setlength{\parskip}{0pt}}
\setcounter{secnumdepth}{-2}

%% ER Settings

%% Page style, header and footer
\usepackage[headsepline]{scrlayer-scrpage}
\pagestyle{scrheadings}
\lohead{\textsc{Heinig}}
\rohead{\textit{Entangled Religions} 14.1 (2023)}
%\cofoot[]{\thepage}
\cofoot[]{}

%% Page margin

% We don’t use the page foot except on the first page.
% This leads to a large bottom margin on most pages.
% To solve this, we’re setting the bottom margin manually
% and set footskip to 0pt. This completely merges the
% foot into the text area.
% We’re then using \enlargethispage{-5\baselineskip} to
% create space for the foot on the first page.
\usepackage[bottom=2.8cm,footskip=0pt]{geometry}

%% Force loading of required packages
\usepackage{booktabs,graphicx}

%% Bookmarks/Outline
\usepackage{bookmark}

%% Text settings
\tolerance=1000
\widowpenalty=10000
\clubpenalty=10000

%% Colors
\definecolor{parnum}{gray}{0.5}

%% Figure captions
\KOMAoptions{captions=tableabove}
\addtokomafont{caption}{\normalfont\small}
\addtokomafont{captionlabel}{\bfseries}
\renewcommand*{\captionformat}{\quad}

%% page headings
\setkomafont{pageheadfoot}{%
\normalfont\normalcolor\small}
\setkomafont{author}{\large\scshape}
\setkomafont{date}{\large}

%% languages/scripts
\newfontfamily\hebrewfont[Script=Hebrew]{FreeSerif}
\newfontfamily\arabicfont[Script=Arabic]{Scheherazade}
\newfontfamily\syriacfont[Script=Syriac]{Noto Sans Syriac Estrangela}
\newfontfamily\greekfont[Script=Greek]{Gentium Plus}
\newfontfamily\copticfont[Script=Coptic]{Noto Sans Coptic}
\newfontfamily\ethiopicfont[Script=Ethiopic]{Noto Serif Ethiopic}
\newfontfamily\amharicfont[Script=Ethiopic]{Noto Serif Ethiopic}

%% enumerations

\renewcommand{\labelenumi}{\arabic{enumi}.}
\renewcommand{\labelenumii}{\arabic{enumi}.\arabic{enumii}}
\renewcommand{\labelenumiii}{\arabic{enumi}.\arabic{enumii}.\arabic{enumiii}}

%% figure positioning
\makeatletter
\def\fps@figure{tbp}
\makeatother

%% set floats above footnotes
\usepackage[bottom]{footmisc}

%% End ER Settings


%% pandoc-tablenos: required package
\usepackage{caption}

%% pandoc-tablenos: environment to disable table caption prefixes
\makeatletter
\newcounter{tableno}
\newenvironment{tablenos:no-prefix-table-caption}{
  \caption@ifcompatibility{}{
    \let\oldthetable\thetable
    \let\oldtheHtable\theHtable
    \renewcommand{\thetable}{tableno:\thetableno}
    \renewcommand{\theHtable}{tableno:\thetableno}
    \stepcounter{tableno}
    \captionsetup{labelformat=empty}
  }
}{
  \caption@ifcompatibility{}{
    \captionsetup{labelformat=default}
    \let\thetable\oldthetable
    \let\theHtable\oldtheHtable
    \addtocounter{table}{-1}
  }
}
\makeatother
\ifXeTeX
  % Load polyglossia as late as possible: uses bidi with RTL langages (e.g. Hebrew, Arabic)
  \usepackage{polyglossia}
  \setmainlanguage[]{english}
  \setotherlanguage[]{greek}
  \setotherlanguage[]{syriac}
  \setotherlanguage[]{amharic}
  \setotherlanguage[]{russian}
  \setotherlanguage[]{hebrew}
  \setotherlanguage[]{farsi}
  \setotherlanguage[]{coptic}
  \setotherlanguage[]{arabic}
\else
  \usepackage[greek,syriac,amharic,russian,hebrew,farsi,coptic,arabic,main=english]{babel}
% get rid of language-specific shorthands (see #6817):
\let\LanguageShortHands\languageshorthands
\def\languageshorthands#1{}
  \newcommand{\textgreek}[2][]{\foreignlanguage{greek}{#2}}
  \newenvironment{greek}[2][]{\begin{otherlanguage}{greek}}{\end{otherlanguage}}
  \newcommand{\textsyriac}[2][]{\foreignlanguage{syriac}{#2}}
  \newenvironment{syriac}[2][]{\begin{otherlanguage}{syriac}}{\end{otherlanguage}}
  \newcommand{\textamharic}[2][]{\foreignlanguage{amharic}{#2}}
  \newenvironment{amharic}[2][]{\begin{otherlanguage}{amharic}}{\end{otherlanguage}}
  \newcommand{\textrussian}[2][]{\foreignlanguage{russian}{#2}}
  \newenvironment{russian}[2][]{\begin{otherlanguage}{russian}}{\end{otherlanguage}}
  \newcommand{\texthebrew}[2][]{\foreignlanguage{hebrew}{#2}}
  \newenvironment{hebrew}[2][]{\begin{otherlanguage}{hebrew}}{\end{otherlanguage}}
  \newcommand{\textfarsi}[2][]{\foreignlanguage{farsi}{#2}}
  \newenvironment{farsi}[2][]{\begin{otherlanguage}{farsi}}{\end{otherlanguage}}
  \newcommand{\textcoptic}[2][]{\foreignlanguage{coptic}{#2}}
  \newenvironment{coptic}[2][]{\begin{otherlanguage}{coptic}}{\end{otherlanguage}}
  \newcommand{\textarabic}[2][]{\foreignlanguage{arabic}{#2}}
  \newenvironment{arabic}[2][]{\begin{otherlanguage}{arabic}}{\end{otherlanguage}}
\fi
\ifLuaTeX
  \usepackage{selnolig}  % disable illegal ligatures
\fi

%% Footnotes (load after polyglossia!)
\deffootnote{\leftmargin}{1em}{%
\makebox[\leftmargin][l]{\thefootnotemark}%
}

\newlength{\cslhangindent}
\setlength{\cslhangindent}{\leftmargin}
\newlength{\csllabelwidth}
\setlength{\csllabelwidth}{3em}
\newenvironment{CSLReferences}[2] % #1 hanging-ident, #2 entry spacing
 {% don't indent paragraphs
  \setlength{\parindent}{0pt}
  % turn on hanging indent if param 1 is 1
  \ifodd #1 \everypar{\setlength{\hangindent}{\cslhangindent}}\ignorespaces\fi
  % set entry spacing
  \ifnum #2 > 0
  \setlength{\parskip}{#2\baselineskip}
  \fi
 }%
 {}
\usepackage{calc}
\newcommand{\CSLBlock}[1]{#1\hfill\break}
\newcommand{\CSLLeftMargin}[1]{\parbox[t]{\csllabelwidth}{#1}}
\newcommand{\CSLRightInline}[1]{\parbox[t]{\linewidth - \csllabelwidth}{#1}\break}
\newcommand{\CSLIndent}[1]{\hspace{\cslhangindent}#1}

\title{Horrordatei zum Testen vom neuen Textsatz}
\usepackage{etoolbox}
\makeatletter
\providecommand{\subtitle}[1]{% add subtitle to \maketitle
  \apptocmd{\@title}{\par {\large #1 \par}}{}{}
}
\makeatother
\subtitle{Untertitel zur Horrordatei}
\author{Julia Heinig\\Centrum für Religionswissenschaftliche Studien,
Ruhr-Universität Bochum}
\date{}

\begin{document}
%% No additional space between sentences
\frenchspacing

\cleardoubleoddpage
\thispagestyle{empty}
\begin{singlespace}

%% Logo
\hfill
\includegraphics[height=3\baselineskip]{Logo_ER.pdf}
\noindent\rule{\textwidth}{\heavyrulewidth}

%% Head
\vspace*{.2\baselineskip}
\noindent\begin{minipage}[c]{.13\textwidth}
\scriptsize
\parbox{.2\textwidth}{
\includegraphics[height=2\baselineskip,width=\csname Gin@nat@width\endcsname,keepaspectratio]{cc-by.pdf}}
\end{minipage}
\begin{minipage}[c]{.4\textwidth}
\scriptsize
\raggedright
DOI:~\href{https://doi.org/http://doi.org./10.13154/er.volume.year.ojsID}{http://doi.org./10.13154/er.volume.year.ojsID}\\
License:~\href{http://creativecommons.org/licenses/by/4.0/}{Creative
Commons Attribution 4.0}
\end{minipage}\hfill
\begin{minipage}[c]{.4\textwidth}
\scriptsize
\raggedleft
\textit{Entangled Religions} 14.1 (2023)\\
\href{https://er.ceres.rub.de/}{er.ceres.rub.de}
\end{minipage}

%% Title
\vspace*{2\baselineskip}
{\usekomafont{title}\LARGE\noindent\raggedright Horrordatei zum Testen
vom neuen Textsatz\par}
\vspace*{.5\baselineskip}
{\usekomafont{subtitle}\Large\noindent\raggedright Untertitel zur
Horrordatei\par}
\vspace{2\baselineskip}
{\usekomafont{author}\noindent Julia
Heinig\hspace{.25em}\href{https://orcid.org/0000-0001-8690-0832}{
\includegraphics[height=1em]{orcid_iD.pdf}}\par}
{\usekomafont{author}\normalfont\itshape\normalsize\noindent Centrum für
Religionswissenschaftliche Studien, Ruhr-Universität Bochum\par}
\vspace{\baselineskip}
\vspace{2\baselineskip}
\end{singlespace}
\begin{quote}
\small
\begin{tabular}{@{}p{\linewidth}@{}}
\toprule
{\bfseries\scshape abstract}\hspace{1em}Lorem ipsum dolor sit amet,
consetetur sadipscing elitr, sed diam nonumy eirmod tempor invidunt ut
labore et dolore magna aliquyam erat, sed diam voluptua. At vero eos et
accusam et justo duo dolores et ea rebum. Stet clita kasd gubergren, no
sea takimata sanctus est Lorem ipsum dolor sit amet. Lorem ipsum dolor
sit amet, consetetur sadipscing elitr, sed diam nonumy eirmod tempor
invidunt ut labore et dolore magna aliquyam erat, sed diam voluptua. At
vero eos et accusam et justo duo dolores et ea rebum. Stet clita kasd
gubergren, no sea takimata sanctus est Lorem ipsum dolor sit amet.\\
\midrule[\heavyrulewidth]
{\bfseries\scshape keywords}\hspace{1em}Entangled Religions, religious
contact, Markdown\\
\bottomrule
\end{tabular}
\end{quote}
\vspace{1.8181\baselineskip} % Since we altered line spacing, we need to adjust the vspace here.

\setstretch{1.1}
\hypertarget{text-farbig-rezania_azar}{%
\section{Text farbig (Rezania\_Azar)}\label{text-farbig-rezania_azar}}

\begin{quote}
\leavevmode\marginpar{\textcolor{parnum}{{[}1{]}}}\textbf{Y.9.1}

\leavevmode\marginpar{\textcolor{parnum}{{[}2{]}}}{\emph{hāuuanīm ā
ratūm ā haomō upāit̰ zaraθuštrǝm}} * \emph{pad hāwan radīh}
{{[}{[}\emph{pad hāwan gāh}{]}{]}} \emph{hōm abar raft ō Zardušt}
\end{quote}

\begin{quote}
\leavevmode\marginpar{\textcolor{parnum}{{[}3{]}}}* {{[}\emph{ātrǝm
pairi yaoždaθǝṇtǝm gāθåsca srāuuaiiaṇtǝm}{]}} * \emph{pad ātaxš-gāh
pērāmōn yōǰdahrēnišnīh ka-š} {{[}{[}\emph{ān Ašem-wohū sē}{]}{]}}
\emph{guft} {{[}{[}\emph{kē Frawarānēy ō pēš}{]}{]}}
\end{quote}

\begin{quote}
\leavevmode\marginpar{\textcolor{parnum}{{[}4{]}}}* {\emph{ādim pǝrǝsat̰
zaraθuštrō kō narǝ ahī}} * *u-š az ōy pursīd Zardušt kū kē mard hē
{{[}{[}hād nē pad yašt ī fradom būd az pēš paydāg. u-š dānist kū hōm ōh
rasēd ud ka mad būd ā-š pursīd abāyist mad miθrō upāit zardušt ān paydāg
kū-š šnāxt ēd rāy čē ān zamān abāg yazdān wēš būd ēstād u-š yād āšnāgtar
būd hēnd. u-š ēn fragard warm būd u-š abāyist rāy abāg hōm ul guft.* *
\emph{ast kē ēdōn gōwēd hād} * \emph{Ohrmazd guft ēstād kū harw dō ōh
rasēnd ud ka hōm mad būd ā-š madan šnāsēd.}{]}{]}}
\end{quote}

\begin{quote}
\leavevmode\marginpar{\textcolor{parnum}{{[}5{]}}}* {\emph{yim azǝm
vīspahe aŋhǝ̄uš astuuatō sraēštǝm dādarǝsa xᵛahe gaiiehe xᵛanuuatō
amǝṣ̌ahe}} * \emph{kē man az harwisp axw ī astōmand ā-m nēktar dīd hē
čē-t ān ī xwēš ǰān nēk kard ēstēd ud amarg} {{[}{[}\emph{hād ā-š pad
frārōnīh ā amarg kerd estēd nē ēdōn čiyōn awēšān kē gōšt ī ǰam ǰūd u-šān
andar tan amarg kerd estād tā bē az tan harw kas-ēw amarg}
{[}\ldots{]}{]}{]}}
\end{quote}

\hypertarget{sprachen}{%
\section{Sprachen}\label{sprachen}}

\hypertarget{griechisch}{%
\subsection{Griechisch}\label{griechisch}}

\leavevmode\marginpar{\textcolor{parnum}{{[}6{]}}}\textgreek{ὁποιόν τι
καὶ πρὸς τὴν Μάρθαν εἰπεν ὁ σωτὴρ ἀσχολουμένην περὶ πολλὰ καὶ
περιελκομένην καὶ ταρασσομένην διακονικῶς, τὴν δὲ ἀδελφὴν αἰτιωμένην ὅτι
τὸ ὑπερετεῖν ἀπολιποῦσα τοῖς ποσὶν αὐτοῦ παρακάθηται μαθητικὴν ἄγουσα
σχολήν · ``σὺ περὶ πολλὰ ταράσσῃ·Μαρία δὲ τὴν ἀγαθὴν μερίδα ἐξελέξατο,
καὶ οὐκ ἀφαιρεθήσεται αὐτης.'' οὑτως καὶ τοῦτον ἐκέλυε τῆς
πολυπραγμοσύνης ἀφέμενον ἑνὶ προστετηκέναι καὶ προσκαθέζεσθαι, τῇ χάριτι
τοῦ ζωήν αἰώνιον προστιθέντος}

\hypertarget{chinesich-mukai}{%
\subsection{Chinesich (Mukai)}\label{chinesich-mukai}}

\leavevmode\marginpar{\textcolor{parnum}{{[}7{]}}}This paper discusses
the complexity of the religious traditions in Quanzhou (Fujian, China),
the largest international trade port under Mongol rule. The contribution
of presumed Persian Muslim Pu Shougeng 蒲壽庚to the reconstruction of a
Taoist-Buddhist shrine was taken as the case study. The external
conditions surrounding his composite religious act (beyond private
beliefs) were also observed in terms of individual goals, backgrounds,
and social networks. For this purpose, the author presents the Chinese
stone inscription from Quanzhou (in Fujian, China) titled ``Zhong jian
Qingyuan Chunyang dong ji 重建清源純陽洞記 (Record of Reconstruction of
the Chunyang Cave in Qingyuan Mountain),'' dated to the fourth year of
Hou-Zhiyuan 後至元 (1338) during the Yuan period.

\hypertarget{syriac-mustataprobleme}{%
\subsection[Syriac (Mustata)]{\texorpdfstring{Syriac
(Mustata\footnote{Hier gab es besonders Probleme mit den Fußnoten
  innerhalb des Syrischen --\textgreater{} falsche Schriftart in der
  Ausgabe vs Durcheinanderbringen der Reihenfolge der Textauszüge, wenn
  FN aus Sprachauszeichnungen ausgenommen werden.})}{Syriac (Mustata)}}\label{syriac-mustataprobleme}}

\begin{itemize}
\item
  \leavevmode\marginpar{\textcolor{parnum}{{[}8{]}}}\textbf{fol.~33v-58r}:
  Kadavil Chandy's \emph{mēmrā} on the Eucharist. The title and the tune
  anthem is missing; it begins with the first verse under the heading
  \emph{ālap} (\textsyriac{ܐ}) of the acrostic. Incipit: (\textsyriac{܀
  ܐܵܒܵܐ ܠܲܒܪܹܗ̣ ܝܼܚܝܼܕܵܝܵܐ ܠܥܵܠܲܡ ܫܲܕܲܪ܀ ܕܓܹܢܣܵܐ ܕܐܢ̈ܫܵܐ ܕܲܒܝܼܕ ܗܘܵܐ
  ܢܒܲܫܩܲܪ{[}\^{}9{]} ܐܵܦܿ ܢܹܣܥܘܿܪ܀}). The scribe copied the tune anthem,
  which usually follows the title of the poem, only at the end of the
  \emph{mēmrā} (on fol.~57v-58r).
\item
  \leavevmode\marginpar{\textcolor{parnum}{{[}9{]}}}\textbf{fol.~58r-69v}:
  Kadavil Chandy's \emph{mēmrā} on the Syriac language. Title and
  incipit: (\textsyriac{ܫܝܼܪ̈ܬܵܐ ܕܥܲܠ ܩܘܼܠܵܣ ܠܹܫܵܢܵܐ ܣܘܼܪܝܝܵܐ ܒܿܩܲܠ
  ܣܓܖܝܼܢܲܢ{[}\^{}10{]} ܒܲܪܟܝܼܢܲܢ. ܠܚܲܕ ܟܝܵܢܵܐ ܡܥܠܝܐ܀ ܡܲܘܕܹܝܢܲܢ. ܒܲܫܡܹܗ
  ܡܬܘܿܡܵܝܵܐ܀ ܕܲܐܒܵܐ ܘܲܒܼܪܵܐ. ܘܪܘܼܚܵܐ ܕܩܘܼܕܫܐ.}).
\item
  \leavevmode\marginpar{\textcolor{parnum}{{[}10{]}}}\textbf{fol.~69v-71r}:
  Canticle of glorification in praise of St.~Mary. Title and incipit:
  (\textsyriac{ܬܫܒܘܿܚܬܵܐ ܕܡܲܪܬܿܝ ܡܲܪܝܲܡ{[}\^{}10a{]} ܒܿܪܝܼܟܼ ܚܢܵܢܵܐ
  ܕܲܒܼܛܲܝܒܘܼܬܹܗ̣ ܦܲܪܢܹܣ ܚܲܝܝ̈ܢ ܒܲܢܒܝܼܘܼܬܼܵܐ܀}).
\item
  \leavevmode\marginpar{\textcolor{parnum}{{[}11{]}}}\textbf{fol.~71r-72r}:
  Intercession before the altar. Title and incipit: (\textsyriac{ܒܥܘܬܐ
  ܖܩܖܡ ܡܖܒܚܐ{[}\^{}10b{]} ܐܲܒܵܐ ܕܩܘܼܫܬܐ ܗܐ ܒܪܵܟ ܕܹܒܚܵܐ ܕܲܡܪܥܹܐ ܠܵܟܼ܀}).
\item
  \leavevmode\marginpar{\textcolor{parnum}{{[}12{]}}}\textbf{fol.~72v-74r}:
  Untitled hymn about the Crucifixion of Christ. In the first part of
  the hymn the Greek \emph{Trìs hágion} hymn is meant to be repeated
  after some of the stanzas, while after other stanzas the words of
  Christ ``My people, what have I done to thee?'' are to be repeated.
  Incipit: (\textsyriac{ܥܲܡܝ ܡܵܢܵܐ ܥܹܒܼܕܲܬܼ{[}\^{}11{]} ܠܵܟܼ܆ ܥܢܝܼܢܝ
  ܒܿܡܲܢ ܐܲܟܼܪܝܼܬܵܟܼ܀ ܕܡ̣ܢ ܐܲܪܥܵܐ ܕܡܹܨܪܹܝܢ: ܐܲܦܿܩܹܬܼ ܠܲܓܦܹܬܵܟ܀}).
\end{itemize}

\hypertarget{geez-gebreyes-auch-miszelle-von-dege-muxfcller}{%
\subsection{Ge'ez (Gebreyes, auch: Miszelle von
Dege-Müller)}\label{geez-gebreyes-auch-miszelle-von-dege-muxfcller}}

\leavevmode\marginpar{\textcolor{parnum}{{[}13{]}}}It is obvious that
the edict was strictly obeyed among the Christian subjects of the
kingdom. But for those who disobeyed, the punishment was also specified:
\textamharic{ወዘንተ፡ ዘኢገብረ፡ ይትበርበር፡ ቤቱ፡ ወይትኰነን፡ ሥጋሁ፡ እንዘ፡ ይብል፡ አዘዘ፡ ንጉሥ፡
በውስተ፡ ኵሉ፡ ዓለም፡} (`He who did not heed {[}this prescription{]} had his
house pillaged and corporal punishment was also administered to the
guilty party to recall that the king should be obeyed by everybody')

\hypertarget{russich-lecke}{%
\subsection{Russich (Lecke)}\label{russich-lecke}}

\leavevmode\marginpar{\textcolor{parnum}{{[}14{]}}}\textrussian{«Мне
кажется, что и в советское время люди чувствовали негласную иерархию.
Русские были титульной нацией. Белорусы и украинцы --- младшими
братьями. Потом кавказские народы --- они не славяне, от нас дальше, но
у них древняя культура, солнце, море и вкусная еда. И, наконец, народы
Средней Азии, чьи территории были насильственно присоединены в ходе
колониальных войн, а в советская время достаточно произвольно поделены
на республики», --- размышляет Виктория.}

\hypertarget{hebruxe4isch-haim}{%
\subsection{Hebräisch (Haim)}\label{hebruxe4isch-haim}}

\leavevmode\marginpar{\textcolor{parnum}{{[}15{]}}}\texthebrew{צׄונין
גופת י'י כודאה גׄאדה ריכתה אמדן תני תו ואשכארה כרדה אמד שַרְמְגָּאה תו
פא} \texthebrew{והאריהא תו אבא דוסתאן תו ואבא המא בותיהא זישתיהא תו
וצׄון כוניהא פסראן תו אן כי בי דאדי פא אישאן} (\emph{čunīn guft
\textbf{adonay} ḫ\textsuperscript{w}adāh ǧāda-yi riḫta āmadan-i tan-i tu
wa-āškāra karda āmad šarmgāh-i tu pa wahārīhā-yi tu abā dūstān-i tu
wa-abā hamā butīhā-yi zištīhā-yi tu wa-čun ḫūnīha-yi pusarān-i tu ān ki
bi-dādī pa īšān})

\hypertarget{perisch-paragraphnummer-erscheint-in-der-ausgabe-bisher-auf-der-linken-seite}{%
\subsection{Perisch (Paragraphnummer erscheint in der Ausgabe bisher auf
der linken
Seite)}\label{perisch-paragraphnummer-erscheint-in-der-ausgabe-bisher-auf-der-linken-seite}}

\hypertarget{bertolacci}{%
\subsubsection{Bertolacci}\label{bertolacci}}

\begin{farsi}

\begin{quote}
\leavevmode\marginpar{\textcolor{parnum}{{[}16{]}}}\textfarsi{فرغ من
تسوید هذه المجلّدة الداعي لصاحبها أینما کان محمود بن عليّ بن محمّد
\textbar{} ابن عليّ الوندکلي وهي قریة من إحدی قری قاسان حماها الله من
طوارق \textbar{} الحدثان في أوائل ربیع الأوّل من شهور سنة ثمان عشر
وسبعمائة \textbar{} وقد کتب من نسخة سقیمة کثیرة التصحیفات قلیلة
التصحیحات \textbar{} والکاتب في أیدي الزمان أسیر و في قید الهوان کسیر
\textbar{} ومع ذلك لا یقدر علی إدراك معانیها وإبدال ألفاظها وقد
\textbar{} استمسك بالعروة الوثقی أعني ألطاف المولی أن یبلّ (؟) عليّ
\textbar{} ذیل الإغماض وإلّا فاقض ما أنت قاض فإنّه قد بلغت نفسي
\textbar{} بأقصی غایة جهدها ولا یکلّف الله نفساً إلّا وسعها \textbar{}
وصلّی الله علی خیر الأخیار محمّد وصحبه الأبرار.}
\end{quote}

\end{farsi}

\begin{farsi}

\begin{quote}
\leavevmode\marginpar{\textcolor{parnum}{{[}17{]}}}\textfarsi{هو\textbar{}
إنّ هذا الکتاب {کان في مطالعة جدّي} وسیّدي وإسنادي {صدر} الحکماء بدر
العلماء شمس السماء \textbar{} قمر الخضراء النیّر الأعظم ومنوّر العالم
{صدر} الحقیقة {محمّد} الماضي علیه من الله السلام {ثمّ انتقل منه إلی
ولده} \textbar{} العلّامة {وهو أبي} وسيّدي وأستادي فخر آبائي وأجدادي
إمام الحکمة {غیاث} النفوس \textbar{} کاشف الغمّة صاحب الهمّة الإمام
الجامع الغالب علی الشیخ الرئیس والحکیم العظیم الفائق علی \textbar{}
أرسطاطالیس أکمل أهل النظر أستاد البشر العقل الحادي عشر أعني الحضرة
العلیّة البهیّة \textbar{} والسدّة السنیّة الجلیّة الفیلسوفیة {غیاث}
الأنام {المنصور} کاسمه ناصر الشریعة والإسلام \textbar{} {ثمّ انتقل منه
إلی ابنه} وتلمیذه بل أقلّ عبد من عبیده الباسط ذراعیه بالوصید \textbar{}
{محمّد بن علي} (؟) {منصور بن محمّد الحسیني} المشتهر {بصدر الثاني} شرح
الله صدره ورفع قدره.}
\end{quote}

\end{farsi}

\hypertarget{rezania-azva-zweifarbig-markiert}{%
\subsubsection{Rezania Azva (zweifarbig
markiert)}\label{rezania-azva-zweifarbig-markiert}}

\begin{farsi}

\begin{quote}
\begin{enumerate}
\def\labelenumi{(\arabic{enumi})}
\setcounter{enumi}{39}
\tightlist
\item
  \leavevmode\marginpar{\textcolor{parnum}{{[}18{]}}}{زابهاش هروند اسپ و
  فه هروار دم له چایند} * فروزهاش بسیار است و بشمار در نیایند
\end{enumerate}
\end{quote}

\begin{quote}
\begin{enumerate}
\def\labelenumi{(\arabic{enumi})}
\setcounter{enumi}{40}
\tightlist
\item
  \leavevmode\marginpar{\textcolor{parnum}{{[}19{]}}}{فر هوشام لی هروار
  خمیده هزهیشام خر میم بهنام اسپ که سارام اورنگرامان و حمیدگان نوند شکار
  هواند} * فرشتگان بی‌شمار آفریده از ایشان نخستین خرد نخست است که همه
  خردها و آفریدگان زیردست اویند
\end{enumerate}
\end{quote}

\begin{quote}
\begin{enumerate}
\def\labelenumi{(\arabic{enumi})}
\setcounter{enumi}{41}
\tightlist
\item
  \leavevmode\marginpar{\textcolor{parnum}{{[}20{]}}}{نزم مانیستار که
  فدسمند اسپ ویودسرو نویشرامان است} * پس روان سپهر برتر که بس بزرگست و
  سالار همه روانهاست
\end{enumerate}
\end{quote}

\begin{quote}
\begin{enumerate}
\def\labelenumi{(\arabic{enumi})}
\setcounter{enumi}{42}
\tightlist
\item
  \leavevmode\marginpar{\textcolor{parnum}{{[}21{]}}}{یرسریربد وهو فرسار
  سروسریرام اسپ} * پس تنبد و او سالار همه تنهاست و تن‌بد نام سپهر {و
  تن‌بد نام سپهر برترین است} برترین است
\end{enumerate}
\end{quote}

\begin{quote}
\begin{enumerate}
\def\labelenumi{(\arabic{enumi})}
\setcounter{enumi}{42}
\tightlist
\item
  \leavevmode\marginpar{\textcolor{parnum}{{[}22{]}}}{یرسریربد وهو فرسار
  سروسریرام اسپ} * پس تنبد و او سالار همه تنهاست و تن‌بد نام سپهر برترین
  است
\end{enumerate}
\end{quote}

\begin{quote}
\begin{enumerate}
\def\labelenumi{(\arabic{enumi})}
\setcounter{enumi}{43}
\tightlist
\item
  \leavevmode\marginpar{\textcolor{parnum}{{[}23{]}}}{سیامکان و هرنامگان
  و هرنامگانیان و شاورام و تاو رام سارام خمیده هواند هزهو فستام پم
  هیشام} * آزادان و وارستگان و تنان و تنانیان و گوهرها و ناگوهرها همه
  آفریده اویند از او آفرین برایشان
\end{enumerate}
\end{quote}

\begin{quote}
\leavevmode\marginpar{\textcolor{parnum}{{[}24{]}}}{[}\ldots{]}
\end{quote}

\begin{quote}
\begin{enumerate}
\def\labelenumi{(\arabic{enumi})}
\setcounter{enumi}{46}
\tightlist
\item
  \leavevmode\marginpar{\textcolor{parnum}{{[}25{]}}}{فه سام زمرپان وای}
  * بنام مهربان خدای
\end{enumerate}
\end{quote}

\begin{quote}
\begin{enumerate}
\def\labelenumi{(\arabic{enumi})}
\setcounter{enumi}{47}
\tightlist
\item
  \leavevmode\marginpar{\textcolor{parnum}{{[}26{]}}}{تشتاریدن رام برج
  ادهابیغی و سابیغی} * پاک شدن دو گونه است امیغی و رواییی
\end{enumerate}
\end{quote}

\begin{quote}
\begin{enumerate}
\def\labelenumi{(\arabic{enumi})}
\setcounter{enumi}{48}
\tightlist
\item
  \leavevmode\marginpar{\textcolor{parnum}{{[}27{]}}}{هابیغی منا درافه
  افرکنون له بر تن و هموزیدگیها سجردن} * امیغی دل را به بدی نه بستن و
  نکوهیدگیها ستردن {مانند خشم و کام از دل زدودن}
\end{enumerate}
\end{quote}

\begin{quote}
\begin{enumerate}
\def\labelenumi{(\arabic{enumi})}
\setcounter{enumi}{49}
\tightlist
\item
  \leavevmode\marginpar{\textcolor{parnum}{{[}28{]}}}{و سابیغی هانچیم دم
  پرکتار یاج هاسد سلودن} * و رواییی آنچه در آشکار بد باشد زدودن {چون
  اویژگی و ناپازی آشکاری}
\end{enumerate}
\end{quote}

\begin{quote}
\begin{enumerate}
\def\labelenumi{(\arabic{enumi})}
\setcounter{enumi}{50}
\tightlist
\item
  \leavevmode\marginpar{\textcolor{parnum}{{[}29{]}}}{وهیم تشتاریدن فه
  جریفترپامد} * و این پاک شدن به آب یفتر باشد {و یفتر آبیست که رنگ و بوی
  و مزه او نگشته بود و بدبوی نشده ورنه گلاب و مانند آن پاک‌ون ستوده جم
  است.}
\end{enumerate}
\end{quote}

\begin{quote}
\begin{enumerate}
\def\labelenumi{(\arabic{enumi})}
\setcounter{enumi}{51}
\tightlist
\item
  \leavevmode\marginpar{\textcolor{parnum}{{[}30{]}}}{وجرسود دم کاو
  سمید} * و آب کرد در خورد تنه و توش {باید دانست که آب کرد آنرا گویند که
  تن و چیز بدان پاک شود و آن در خورد تن آمده پس در خورد تن پیل رودی و
  مردم را آنمایه که درو سراپا فرو شود و بهر پشه همینه}
\end{enumerate}
\end{quote}

\end{farsi}

\hypertarget{koptisch-colditz}{%
\subsection{Koptisch (Colditz)}\label{koptisch-colditz}}

\leavevmode\marginpar{\textcolor{parnum}{{[}31{]}}}BeDuhn and Gardner
have argued on the basis of \emph{Dublin Kephalaia} 353.28
{[}\textcoptic{ⲥ}{]}\textcoptic{ⲁⲡⲱⲣⲏⲥ ⲡ}{[}\textcoptic{ⲣ̅ⲣⲟ
ⲛ̅ⲧⲟ}{]}\textcoptic{ⲩⲣⲁⲛ} /\emph{{[}S{]}apōrēs p-{[}erro en-to{]}uran}/
``{[}Š{]}ābuhr the {[}king of T{]}ūrān'' that this king was Šābuhr, son
of the Sasanian king of kings Šābuhr I (240--272 C.E.), who is mentioned
in the inscription of the latter at the Kaʿbe-ye Zardošt at Naqš-e
Rostam (ca. 260--262 C.E.) as ``king of the Mesene'' {[}ŠKZ § 34: MP.,
Parth. \emph{mēšān-šah}, Gr. \textgreek{Μησανηνῶν βασιλέως}.

\hypertarget{hieroglyphen-shalomi-hen}{%
\subsection{Hieroglyphen (Shalomi-Hen)}\label{hieroglyphen-shalomi-hen}}

\leavevmode\marginpar{\textcolor{parnum}{{[}32{]}}}The attestation of
\includegraphics{figure1.svg} -- \emph{I͗hwi͗w} Ihuiu/\emph{ᵓa-hu-ᵓu}
comes from the Biography of \emph{H̱ty}, an overseer of the quarries,
dated to the Eleventh Dynasty (ca. 2000 BCE).{[}\^{}3{]} Although this
attestation is too early for the time frame set for this paper, it is
discussed here because some scholars consider
\includegraphics{figure1.svg} -- \emph{I͗hwi͗w} Ihuiu/\emph{ᵓa-hu-ᵓu} to
be an early orthography of the later YHWA/YHA/YH.

\hypertarget{arabisch}{%
\subsection{Arabisch}\label{arabisch}}

\begin{Arabic}

\leavevmode\marginpar{\textcolor{parnum}{{[}33{]}}}(3) \ldots{} ثم قالوا
\textbf{كي ايل ديعو{[}ث{]}} (4) \textbf{יי'} وفيه قولان الواحد \ldots{}
والاخر ان الله عارف الخفيات عن الناس فمن تعاطا الكلام (6) عليها فقد تكلم
بما ينكره الله عليه وفي مثل ذلك قال داود عليه السلم (7) \textbf{ولوا
هلاختي بغذولوث وبنيفلاوث مماني}

\end{Arabic}

\hypertarget{tabellen-und-bilder}{%
\section{Tabellen und Bilder}\label{tabellen-und-bilder}}

\hypertarget{kruxfcgerradermacher-miszelle}{%
\subsection{Krüger/Radermacher
Miszelle}\label{kruxfcgerradermacher-miszelle}}

\leavevmode\marginpar{\textcolor{parnum}{{[}34{]}}}See table
\ref{tbl:table1} and figure \ref{fig:figure1}:

\begin{longtable}[]{@{}
  >{\raggedright\arraybackslash}p{(\columnwidth - 6\tabcolsep) * \real{0.22}}
  >{\raggedright\arraybackslash}p{(\columnwidth - 6\tabcolsep) * \real{0.39}}
  >{\raggedright\arraybackslash}p{(\columnwidth - 6\tabcolsep) * \real{0.31}}
  >{\raggedright\arraybackslash}p{(\columnwidth - 6\tabcolsep) * \real{0.09}}@{}}
\caption{Inventory of the collection. \label{tbl:table1}}\tabularnewline
\toprule
\textbf{Object} & \textbf{Designation according to Worms (German)} &
\textbf{Designation according to the Inventory (German)} &
\textbf{Inv.-No.} \\
\midrule
\endfirsthead
\toprule
\textbf{Object} & \textbf{Designation according to Worms (German)} &
\textbf{Designation according to the Inventory (German)} &
\textbf{Inv.-No.} \\
\midrule
\endhead
Conical stone & Phallus-Stein & Phallus-Stein, Australien & 12 \\
Clapstick & eingravierte Holzstäbe; Klabgstäbe; Xylophone & Klangholz,
Australien & 13 \\
Clapstick & & Klangholz, Australien & 14 \\
Spear point & Speerspitzen & Speerspitze, Australien & 15 \\
Spear point & & Speerspitze, Australien & 16 \\
lancet-shaped mother-of-pearl plates & mythologische Perlmutter-blätter;
lanzetförmige Perl-mutterstücke & Lanzette, Perlmutter, Australien &
17 \\
lancet-shaped mother-of-pearl plates & & Lanzette, Perlmutter,
Australien & 18 \\
Stone tools (12) & Klein-Steinwerkzeuge & Steine, Australien {[}partly
lost{]} & 18,1--6 \\
Headdress & Kopfschmuck & Schleier & 19 \\
Stone tools (50) & 50 Steinwerkzeuge & Sammlung von Steinwerkzeugen
{[}partly lost{]} & 20 \\
Thread cross & Fadenkreuz & Fadenkreuz vom Grabkult der Australschwarzen
& 23 \\
Thread cross & & {[}not inventoried/lost{]} & \\
Tjurunga (stone) & tellergrosses Stein-Tjuringa & {[}not
inventoried/lost{]} & \\
Tjurunga (wood) & großes Tjuringa (90 cm) & {[}not inventoried/lost{]}
& \\
Tjurunga (wood) & zwei kleinere Tjuringa & {[}later inventoried{]} &
14,2? \\
Tjurunga (wood) & & {[}later inventoried{]} & 14,3? \\
Lara (wood) & eingraviertes Holz, Tjuringa ähnlich & {[}later
inventoried{]} & 14,1? \\
Tjurunga (wood)

{[}?{]} & Tjuringa-Schwirrholz {[}Antweiler{]} {[}?{]} & {[}not
inventoried/lost{]} {[}?{]} & \\
Stone axe & Steinbeil, gehämmert und geschliffen & {[}not
inventoried/lost{]} & \\
Stone knife & Steinmesser für Initiationszeremonien & {[}not
inventoried/lost{]} & \\
Pastel drawings (5) & 5 Pastellzeichnungen eines australischen
Eingeborenen & {[}not inventoried/lost{]} & \\
\bottomrule
\end{longtable}

\begin{figure}
\hypertarget{fig:figure1}{%
\centering
\includegraphics[width=1\textwidth,height=\textheight]{figure1.png}
\caption[Conical Stone, Religionskundliche Sammlung der Westfälischen
Wilhelms-Universität Münster, Inv. No.~12, © Institut für
Religionswissenschaft, WWU Münster, 2021)]{Conical Stone,
Religionskundliche Sammlung der Westfälischen Wilhelms-Universität
Münster, Inv. No.~12, © Institut für Religionswissenschaft, WWU Münster,
2021)\footnotemark{}}\label{fig:figure1}
}
\end{figure}
\footnotetext{Some of the objects examined in this article are subject
  to secrecy in their societies of origin and are considered
  secret/sacred. This applies in particular to the tjuringas, which,
  according to tradition, may not be viewed by women and/or outsiders.
  On the other hand, social-scientific research is increasingly required
  to ensure gender-neutral access (see Raabe 2017). Faced with this
  conflict of accessibility and restrictions, the authors have chosen to
  depict the objects discussed here through simplified line drawings
  that are not equivalent to original photographs.}

\hypertarget{stadlbauer}{%
\subsection{Stadlbauer}\label{stadlbauer}}

\begin{longtable}[]{@{}lllll@{}}
\caption{Asylum applications from the European Union by country of
origin in the years 2014 and 2015 {[}translated from Bundesamt 2016{]}.
\label{tbl:asylum}}\tabularnewline
\toprule
Rank & State of Origin & 2014 & 2015 & Shift 2014/2015 \\
\midrule
\endfirsthead
\toprule
Rank & State of Origin & 2014 & 2015 & Shift 2014/2015 \\
\midrule
\endhead
1 & Syria & 122,065 & 368,400 & +201.8\% \\
2 & Afghanistan & 41,405 & 181,360 & +338.0\% \\
3 & Iraq & 21,365 & 124,905 & +484.6\% \\
4 & Cosovo & 37,890 & 72,465 & +91.3\% \\
5 & Albania & 16,950 & 67,740 & +299.6\% \\
6 & Pakistan & 22,220 & 47,840 & +115.3\% \\
7 & Eritrea & 36,945 & 34,105 & -7.7\% \\
8 & Nigeria & 20,065 & 31,165 & +55.3\% \\
9 & Serbia & 30,840 & 30,050 & -2.6\% \\
10 & Iran & 10,905 & 26,550 & +143.5\% \\
\bottomrule
\end{longtable}

\begin{tablenos:no-prefix-table-caption}

\begin{longtable}[]{@{}
  >{\raggedright\arraybackslash}p{(\columnwidth - 4\tabcolsep) * \real{0.06}}
  >{\raggedright\arraybackslash}p{(\columnwidth - 4\tabcolsep) * \real{0.47}}
  >{\raggedright\arraybackslash}p{(\columnwidth - 4\tabcolsep) * \real{0.47}}@{}}
\toprule
\endhead
1 & \textbf{\emph{Bevor war ich ein Fisch in einem kleinen Aquarium.}} &
\textbf{Before, I was a fish in a small aquarium.} \\
2 & \emph{Aber jetzt, nein, ich liebe Gott.} & But now, no, I love
God. \\
3 & \emph{Bevor Gott sagte, du musst das, das, das} & Before, God said
you have to do this, this, this. \\
4 & \emph{Aber jetzt weiß ich, mein Gott ist groß.} & But now I know my
God is large. \\
5 & \emph{Er ist freundlich.} & He is friendly. \\
6 & \emph{Ich habe keine Angst.} & I have no fear. \\
7 & \emph{Aber bevor hatte ich viel Angst.} & But before I had a lot of
fear. \\
8 & \textbf{\emph{Und jetzt bin ich in einem großen Meer.}} &
\textbf{And now I am in a large ocean.} \\
9 & \emph{Viele Möglichkeiten.} & Many opportunities. \\
10 & \emph{Ich kann denken.} & I can think. \\
11 & \emph{Ich kann fragen.} & I can ask. \\
12 & \emph{Ich kann lesen, lernen.} & I can read, learn. \\
13 & \emph{Und ich bin frei.} & And I am free. \\
\bottomrule
\end{longtable}

\end{tablenos:no-prefix-table-caption}

\hypertarget{original-date-in-2-in-text-citations-hintereinander-in-einer-klammer-in-warner}{%
\section{Original Date (in 2 in-text citations hintereinander in einer
Klammer, in
Warner)}\label{original-date-in-2-in-text-citations-hintereinander-in-einer-klammer-in-warner}}

\leavevmode\marginpar{\textcolor{parnum}{{[}35{]}}}Omidsalar
(\protect\hyperlink{ref-omidsalar_ali_2010}{{[}1389{]} 2010--2011},
\protect\hyperlink{ref-omidsalar_justari_2014}{{[}1393{]} 2014--2015})
has forcefully pushed back against scholarly assumptions of Rabīʿ's
total hostility to the \emph{Shāh nāma} and its contents. While I share
Omidsalar's opposition to a simplistic reading of the \emph{ʿAlī nāma}
as evidence for a generalised hostility between Shīʿī Muslims and the
stuff of Iranian `national' epic
(\protect\hyperlink{ref-omidsalar_ali_2010}{{[}1389{]} 2010--2011,
52--54}), I will argue below that the antipathy that Rabīʿ's poem
projects towards Ferdowsi is real and substantial, expressing not a
universal sentiment but the particular mimetic concerns of Rabīʿ
himself. Specific points of divergence with Omidsalar's arguments in
this regard will be addressed below as and when they appear.

\hypertarget{weitere-sonderfuxe4lle}{%
\section{Weitere Sonderfälle}\label{weitere-sonderfuxe4lle}}

\hypertarget{bogdan-2023}{%
\subsection{Bogdan 2023}\label{bogdan-2023}}

\leavevmode\marginpar{\textcolor{parnum}{{[}36{]}}}5°=6\textsuperscript{□}

\leavevmode\marginpar{\textcolor{parnum}{{[}37{]}}}spiritual legitimacy
to found the Order of the Silver Star, or A∴A∴

\hypertarget{manuelle-reference-list}{%
\subsection{Manuelle Reference List}\label{manuelle-reference-list}}

\hypertarget{bibliography}{%
\section*{References}\label{bibliography}}
\addcontentsline{toc}{section}{References}

\hypertarget{refs}{}
\begin{CSLReferences}{1}{0}
\leavevmode\vadjust pre{\hypertarget{ref-omidsalar_ali_2010}{}}%
Omidsalar, Mahmoud. (1389) 2010--2011. {``ʿAlī nāma, Shāh nāma wa naqdī
bar sākhtārhā-yi īdiʾūlūzhīk dar muṭāliʿāt-i ḥamāsī {[}The \emph{ʿAlī
nāma}, the \emph{Shāh nāma} and a Critique of the Ideological Structures
in the Study of Epic{]}.''} \emph{Ḍamīma-yi Āyīna-yi Mīrāth} 20: 35--56.

\leavevmode\vadjust pre{\hypertarget{ref-omidsalar_justari_2014}{}}%
---------. (1393) 2014--2015. {``Justārī dar ʿadam-i taʿāruḍ-i ʿAlī nāma
bā Shāh nāma-yi Firdawsī {[}An Investigation of the Lack of Conflict
Between the \emph{ʿAlī nāma} and the \emph{Shāh nāma} of Ferdowsi{]}.''}
\emph{Adyān, Madhāhib wa ʿIrfān} 2: 11--26.

\end{CSLReferences}

\end{document}
